\chapter{Derivate parțiale}

\section{Funcții de mai multe variabile}

În general, putem lucra cu funcții de forma $ f : \RR^n \to \RR $,
funcții care acceptă $ n $ variabile, deci $ f(x_1, \dots, x_n) = y \in \RR $.

Cazurile pe care le vom folosi cel mai des sînt $ n = 2 $ și $ n = 3 $.

Dacă $ f : \RR^2 \to \RR $ este o funcție de două variabile, în loc
de $ f(x, y) $, putem gîndi că avem o funcție de forma $ f_x(y) $, pentru
orice $ x \in \RR $.

Evident, putem avea inclusiv $ f : \RR^n \to \RR^m $, adică funcții de
mai multe variabile, care returnează mai multe variabile.

%%%%%%%%%%%%%%%%%%%%%%%%%%%%%%%%%%%%%%%%%%%%%%%%%%%%%%%%%%%%%%%%%%%%%%

\section{Operatori diferențiali}

În cele ce urmează, presupunem că lucrăm cu $ f : \RR^3 \to \RR, f = f(x, y, z) $.

Putem defini \emph{derivata parțială} a lui $ f $ după variabila $ x $
ca fiind derivata obținută prin tratarea lui $ y $ și $ z $ drept parametri.
Notația este $ \dfrac{\partial f}{\partial x} $ sau $ f_x $ sau $ \partial_x f $.
\index{derivate parțiale}

Similar se pot defini și celelalte derivate parțiale. De exemplu:
\[
  f(x, y, z) = 3x^2 y + 5xe^y + \sin(yz).
\]
Avem:
\[
  f_x = 6xy + 5e^y, \quad f_y = 3x^2 + 5xe^y + z\cos(yz), \quad %
  f_z = y\cos yz.
\]
\index{derivate parțiale!gradient}
\index{derivate parțiale!divergență}
\index{derivate parțiale!nabla}
\index{derivate parțiale!laplacian}
\index{derivate parțiale!laplacian!funcție armonică}
Pentru asemenea funcții, se definesc următorii \emph{operatori diferențiali}:
\begin{itemize}
\item \emph{gradient}, notat $ \dr{grad}f = \nabla f = (f_x, f_y, f_z) $.
  \emph{Observație:} Gradientul transformă o funcție într-un vector de funcții.
\item \emph{divergență}, notată $ \dr{div} f = \nabla \cdot f = f_x + f_y + f_z $;
\item \emph{laplacianul}, notat $ \Delta f = \dfrac{\partial^2 f}{\partial x^2} + %
  \dfrac{\partial^2 f}{\partial y^2} + \dfrac{\partial^2 f}{\partial z^2} $.

  O funcție se numește \emph{armonică}, dacă $ \Delta f = 0 $.
\end{itemize}

Fiecare dintre acești operatori poate fi înțeles privind \emph{operatorul nabla},
$ \nabla $ ca fiind:
\[
  \nabla = \Big( \frac{\partial}{\partial x}, \frac{\partial}{\partial y}, %
  \frac{\partial}{\partial z} \Big).
\]
Atunci, \emph{gradientul} reprezintă \emph{aplicarea} lui $ \nabla $ pe $ f $,
iar \emph{divergența} reprezintă \emph{produsul scalar} dintre $ \nabla $
și vectorul $ (f, f, f) $. De asemenea, laplacianul poate fi înțeles ca
produsul scalar $ \nabla \cdot \nabla $, aplicat apoi lui $ f $.

%%%%%%%%%%%%%%%%%%%%%%%%%%%%%%%%%%%%%%%%%%%%%%%%%%%%%%%%%%%%%%%%%%%%%%

\section{Diferențiala totală}
\index{derivate parțiale!diferențiala totală}
Dată o funcție de mai multe variabile, pe lîngă derivatele parțiale, se
poate defini și o \emph{diferențială totală}, care conține toate informațiile
din derivatele parțiale.

Pentru o funcție $ f = f(x, y, z) $, aceasta se definește prin:
\[
  df = f_x dx + f_y dy + f_z dz
\]

De asemenea, avem și diferențiala totală de ordinul 2:
\[
  d^2 f = \frac{\partial^2 f}{\partial x^2} dx^2 + \frac{\partial^2 f}{\partial y^2} + %
  2 \frac{\partial^2 f}{\partial x \partial y} dxdy.
\]

%%%%%%%%%%%%%%%%%%%%%%%%%%%%%%%%%%%%%%%%%%%%%%%%%%%%%%%%%%%%%%%%%%%%%%

\section{Exerciții}

1. Pentru următoarele funcții $ f : \RR^3 \to \RR $, calculați
gradientul, divergența și diferențiala totală în punctele
$ (1, 1, 1) $ și $ (1, 2, 3) $:
\begin{enumerate}[(a)]
\item $ f(x, y, z) = 3x^2 y + x\sin z + z\cos y $;
\item $ f(x, y, z) = \sqrt{x^2 + y^2 + z^2} $;
\item $ f(x, y, z) = \ln(xy) + \ln(xz) + \ln(yz) $;
\item $ f(x, y, z) = e^{x + y + z} $;
\item $ f(x, y, z) = \dfrac{x + y}{z} + \dfrac{y + z}{x} + \dfrac{x + z}{y} $;
\item $ f(x, y, z) = \arctan \dfrac{x}{y} + \arctan \dfrac{y}{z} + \arctan \dfrac{z}{x} $.
\end{enumerate}

%%%%%%%%%%%%%%%%%%%%%%%%%%%%%%%%%%%%%%%%%%%%%%%%%%%%%%%%%%%%%%%%%%%%%%

\vspace{1cm}

2. Verificați dacă următoarele funcții sînt armonice:
\begin{enumerate}[(a)]
\item $ f(x, y) = \ln(x^2 + y^2) $;
\item $ f(x, y, z) = \ln(x^2 + y^2 + z^2) $;
\item $ f(x, y) = \sqrt{x^2 + y^2} $;
\item $ f(x, y, z) = \sqrt{x^2 + y^2 + z^2} $;
\item $ f(x, y) = (x^2 + y^2)^{-1} $;
\item $ f(x, y, z) = (x^2 + y^2 + z^2)^{-1} $.
\end{enumerate}




%%% Local Variables:
%%% mode: latex
%%% TeX-master: "../m1-18-19"
%%% End:
